\par\medskip
\noindent
14.3. \underline{Application to an arithmetic characterization of}
$i:\underline M_K\longrightarrow G_K$.

The characterization we have in mind rests on the observation that,
for every prime number $\ell$, the extension (14.2.2) of
$\underline M_K\otimes\mathbb Z_\ell$ by $T_\ell(G_K)$ deduced from
$i$ is known by purely algebraic means.  In particular, the image of
$i$ under the canonical homomorphism

\begin{equation}\tag{14.3.1}
 \operatorname{Hom}(\underline M_K,G_K)\longrightarrow
 \prod_\ell\operatorname{Ext}^1
 \bigl(\underline M_K\otimes\mathbb Z_\ell,T_\ell(G_K)\bigr)
 \qquad .
\end{equation}

\noindent
is known.  In the case (to which one may always reduce by descent) in
which $\underline M_K$ is constant, with value $M$, the preceding
homomorphism may be interpreted as

\begin{equation}\tag{14.3.2}
 \operatorname{Hom}\bigl(M,G_K(K)\bigr)\longrightarrow
 \prod_\ell\operatorname{Hom}
 \bigl(M,H^1(K,T_\ell(G_K))\bigr),
\end{equation}

\noindent
deduced from the canonical homomorphism

\begin{equation}\tag{14.3.3}
 G_K(K)\longrightarrow\prod_\ell H^1(K,T_\ell(G_K))
 \qquad ,
\end{equation}

\noindent
which comes from the connecting homomorphisms in the Kummer exact
sequences

\begin{equation*}
 0\longrightarrow{}_nG_K\longrightarrow G_K
 \xrightarrow{\ n\ }G_K\longrightarrow0\qquad .
\end{equation*}

(The homomorphism (14.3.1) could in fact always be interpreted as the
connecting homomorphism analogous to (14.3.3) for the twisted group
$\check{\underline M}_K\otimes G_K=
\underline{\operatorname{Hom}}(\underline M_K,G_K)$.)  The kernel of
(14.3.3) plainly consists of the elements of $G_K(K)$ which are
\underline{divisible by} $n$ \underline{for every integer}
$n\geq1$, or, as we shall say, which are \underline{infinitely
divisible}.  Thus this already determines $i$ modulo an element of
$\operatorname{Hom}(M,g)$, where $g$ is the subgroup of $G_K(K)$
formed by its infinitely divisible elements:

\begin{equation*}
 g=\bigcap_{n\geq1}nG_K(K)\qquad .
\end{equation*}

We now have the following result.

\medskip
\noindent
\underline{Lemma 14.3.4.}  \underline{Let} $V$ \underline{be a
discrete valuation ring whose residue field} $k$ \underline{has
characteristic} $p>0$ \underline{and is of finite type over its prime
field, and let} $G_K$ \underline{be a group scheme of finite type over}
$K$, \underline{of unipotent rank zero if} $\operatorname{char}K=0$.
\underline{Then the group} $g$ \underline{of infinitely divisible
elements of} $G_K(K)$ \underline{is zero; that is,} (\underline{if}
$G_K$ \underline{is divisible}) \underline{the homomorphism} (14.3.3)
\underline{is injective}.

After making a finite extension $K'/K$ and replacing $V$ by its
normalization in $K'$, we may suppose that $G_K$ admits a composition
series whose successive quotients are of one of the following types:
a constant finite group, an abelian variety, $\mathbb G_{mK}$, or
$\mathbb G_{aK}$.  We are therefore immediately reduced to the case
where $G_K$ itself is of one of these types.  The case of a constant
finite group, or of $\mathbb G_{aK}$, is trivial.  In the abelian case,
$G_K$ may be regarded as the generic fibre of its N\'{e}ron model $G$,
so we are reduced to showing that, if $G$ is a scheme of finite type
over $S$, the subgroup of $G(S)$ formed by the infinitely divisible
points is zero.

This is true of the group $G(k)$, as follows at once from the usual
d\'{e}vissage of $G_o$ and from the \underline{Mordell--Weil theorem}
[16] (which says that $G_o(k)$ is a finitely generated $\mathbb Z$-module
if $G_o$ is an abelian scheme over the field of finite type $k$).
Consequently $g$ is contained in the kernel $G(S)^{(1)}$ of
$G(S)\longrightarrow G(k)$.  The kernel of
$G(S)^{(1)}\longrightarrow G(V/\underline m^{,n+1})$ has a filtration
whose quotients are isomorphic to vector groups over $k$ (after
supposing, as we may, that $S$ is complete), and these are killed by
$p$.  It follows that $g$ is contained in the kernel of the preceding
homomorphism for every $n$, and hence that it is zero.  The case
$G_K=\mathbb G_{mK}$ reduces to the preceding one, since
$G_K(K)=K^*$ is an extension of $\mathbb Z$ (which has no nonzero
infinitely divisible elements) by
$V^*=\mathbb G_m(S)$, and the latter satisfies the same condition by
what we have just proved.

\medskip
\noindent
\underline{Remarks 14.3.4.1.}  Notice that in 14.3.4 it is essential
to require divisibility by \underline{all} integers $\geq1$, and not
only by the powers of a fixed prime $\ell$; equivalently, one must
consider the kernel of (14.3.3), rather than restricting the right-hand
side to a single factor $H^1(K,T_\ell(G_K))$.  Indeed, if $G_K$ comes
from a smooth group scheme $G$ over $S$ and $S$ is henselian, every
element of the kernel of $G(S)\longrightarrow G(k)$ is infinitely
$\ell$-divisible for every $\ell\neq p$!  The argument proving 14.3.4
shows that the subgroup of $G_K(K)$ formed by the infinitely
$p$-divisible elements (that is, the kernel of
$G_K(K)\longrightarrow H^1(K,T_p(G_K))$ when $G_K$ is $p$-divisible)
is a \underline{finite} group.

\medskip
\noindent
14.3.5. The preceding considerations show that, when the residue
field $k$ has characteristic $p>0$ and is of finite type in the
absolute sense, the homomorphism $i$ (14.1.5) is completely determined
by its image under (14.3.1), that is, by the extensions (14.2.2).  To
deduce from this a principle characterizing $i$ in the general case,
with no restriction on $k$, observe that in paragraph 7 we were able to
construct the scheme $G$ over $S$ under conditions considerably more
general than those in which $S$ is a complete trait and one has an
abelian variety with semistable reduction over its field of fractions.

More generally, let $S$ denote the spectrum of a normal noetherian ring
$V$, separated and complete for an $\underline m$-adic topology, and
consider a smooth quasi-projective group scheme $A$ over $S$, with
connected fibres.  Suppose that
$A_o=A\times_S S_o$, where
$S_o=\operatorname{Spec}(V/\underline m)$, is an extension of an
abelian scheme $B_o$ by an isotrivial torus $T_o$, and that $A_U$ is an
abelian scheme, where $U\subset S-S_o$ is a nonempty open subset.  In
7.22 we then constructed a scheme $G$ over $S$, an extension of an
abelian scheme $B$ by a torus $T$, characterized up to unique
isomorphism by the isomorphism (14.1.8)

\begin{equation}\tag{14.3.5.1}
 \hat G\xrightarrow{\ \sim\ }\hat A
\end{equation}

\noindent
on the formal completions along $G_o$ and $A_o$.  For every prime
number $\ell$ one deduces an isomorphism (7.3.5)

\begin{equation}\tag{14.3.5.2}
 T_\ell(G)\xrightarrow{\ \sim\ }T_\ell(A)^f\qquad ,
\end{equation}

\noindent
and hence, on restricting to $U$, the structure of an extension of
pro-$\ell$ Barsotti--Tate groups

\begin{equation}\tag{14.3.5.3}
\begin{tikzcd}[row sep=large,column sep=large]
0 \arrow[r] & T_\ell(G_U) \arrow[r] \arrow[d,equal]
 & T_\ell(A_U) \arrow[r]
 & T_\ell(A_U)/T_\ell(A_U)^f \arrow[r] & 0\quad .\\
& T_\ell(A_U)^f
\end{tikzcd}
\end{equation}

Suppose now that a second group scheme $A'$ over $S$ is given,
satisfying the same hypotheses as $A$, together with a biextension $W$
of $(A,A')$ by $\mathbb G_m$ which over $U$ induces a duality between
the abelian schemes $A_U$ and $A'_U$.  (In fact, by VIII 7.1, giving
$W$ is equivalent to giving such a duality between $A_U$ and $A'_U$.)
Then $W$ defines pairings

\begin{equation}\tag{14.3.5.4}
 T_\ell(A_U)\times T_\ell(A'_U)\longrightarrow\mathbb Z_\ell(1)
 \qquad ,
\end{equation}

\noindent
which give rise to the orthogonality theorem (7.4.8):
$T_\ell(A_U)^f$ is the orthogonal of
$T_\ell(A'_U)^t=T_\ell(T'_U)$, where $T'$ is the toric part of the
Raynaud extension $G'={A'}^{\natural}$ associated with $A'$.  Thus one
again obtains an isomorphism

\begin{equation}\tag{14.3.5.5}
 T_\ell(A_U)/T_\ell(A_U)^f
 \simeq\underline M_U\otimes\mathbb Z_\ell\qquad ,
\end{equation}

\noindent
where $\underline M=D(T')$ is the lattice over $S$ dual to the torus
$T'$.  Consequently (14.3.5.3) and (14.3.5.5) give an extension

\begin{equation}\tag{14.3.5.6}
 0\longrightarrow T_\ell(G_U)\longrightarrow T_\ell(A_U)
 \longrightarrow\underline M_U\otimes\mathbb Z_\ell
 \longrightarrow0\qquad .
\end{equation}

With this understood, we may state the following conjecture.
